% 04_propagation_and_acm.tex
% Chapter 4: Atmospheric Propagation and ACM [EXTRA]

\section{Atmospheric Propagation and ACM}

\vspace{0.3cm}

This chapter describes the atmospheric attenuation, illustrative outage, and DVB-S2 ACM calculations used after the clear-sky budget.

\textit{Model scope: the software is version-locked to the recommendation editions cited in this report, including P.618-13. Some recommendations have newer revisions, but changing versions would require a separate code and regression review. The present results use representative meteorological inputs rather than site-specific ITU digital-map extraction and should not be treated as an operational availability guarantee.}

\subsection[ITU-R Slant-Path Attenuation Suite]{ITU-R Slant-Path Attenuation Suite}

\subsubsection{Slant-Path Loss Mechanisms}
The Earth--space path crosses the troposphere, where gases, cloud liquid water, refractive-index fluctuations, and precipitation affect the received carrier. At Ku-band, rain is the largest time-varying attenuation term in the selected 99.9\% design case.

\subsubsection{Model Scope and Inputs}
The ITU-R procedures are used to calculate excess attenuation beyond the fixed clear-sky allowances. Uplink and downlink attenuation are evaluated separately because their frequencies, elevations, and terminal apertures differ.

\subsubsection{Mathematical Formulation}
The total atmospheric attenuation $A_{\text{total}}(p)$ in dB exceeded for an annual time percentage $p$ ($0.001\% \le p \le 5\%$) is calculated by combining gaseous absorption ($A_{\text{gaseous}}$), cloud attenuation ($A_{\text{cloud}}$), tropospheric scintillation ($A_{\text{scintillation}}$), and rain attenuation ($A_{\text{rain}}$) according to \textit{ITU-R P.618-13} \cite{itu618}:
\begin{equation}
    A_{\text{total}}(p) = A_{\text{gaseous}} + \sqrt{ \left( A_{\text{rain}}(p) + A_{\text{cloud}} \right)^2 + A_{\text{scintillation}}(p)^2 }
\end{equation}

The point rainfall rate $R_{0.01}$ is a representative P.837-based input \cite{itu837}. The same $R_{0.01}=42$~mm/h value is used at both terminals; a site-specific study would use separate mapped values for Rome and Ankara. Rain height follows P.839-4 \cite{itu839}, while the frequency- and polarization-dependent $k$ and $\alpha$ coefficients follow P.838-3 \cite{itu838}. The specific attenuation is:
\begin{equation}
    \gamma_{\text{R}} = k R_{0.01}^\alpha
\end{equation}
The effective path length through the rain cell accounts for spatial localisation:
\begin{equation}
    L_{\text{eff}} = L_{\text{s}} \cdot r_{0.01} \cdot v_{0.01}
\end{equation}
where $r_{0.01}$ is the horizontal path reduction factor and $v_{0.01}$ is the vertical adjustment factor. The attenuation exceeded for $0.01\%$ of the year is $A_{0.01} = \gamma_{\text{R}} L_{\text{eff}}$. For other time percentages $p$:
\begin{equation}
    A_p = A_{0.01} \left( \frac{p}{0.01} \right)^{-\left[0.655 + 0.033\ln p - 0.045\ln A_{0.01} - \beta(1-p)\sin\theta\right]}
\end{equation}

Gaseous attenuation uses dry air and water vapor specific zenith attenuations scaled by the slant-path elevation angle (\textit{ITU-R P.676-12} \cite{itu676}):
\begin{equation}
    A_{\text{gaseous}} = \frac{\gamma_{\text{o}} h_{\text{o}} + \gamma_{\text{w}} h_{\text{w}}}{\sin\theta}
\end{equation}

Cloud liquid water attenuation uses columnar liquid water $M$ and specific attenuation $K_{\text{l}}$ (\textit{ITU-R P.840-8} \cite{itu840}):
\begin{equation}
    A_{\text{cloud}} = \frac{M \cdot K_{\text{l}}}{\sin\theta}
\end{equation}

Scintillation uses the P.618-13 procedure with refractivity-related inputs consistent with P.453-14 \cite{itu618,itu453}:
\begin{equation}
    A_{\text{scintillation}}(p) = a(p) \sigma
\end{equation}

\subsubsection{Operational Impact and Results}
At 99.9\% target availability, the model gives \upRainAtten~dB total uplink attenuation at 14~GHz and \downRainAtten~dB total downlink attenuation at 12~GHz. With the selected interference and dynamic-noise assumptions, the fixed-rate end-to-end margin becomes \marginRain~dB.

\begin{figure}[H]
    \centering
    \begin{subfigure}[b]{0.48\textwidth}
        \centering
        \includegraphics[width=\textwidth]{23_itu_attenuation_vs_availability.png}
        \caption{Slant-path attenuation vs. annual link availability.}
        \label{fig:itu_attenuation_curve}
    \end{subfigure}
    \hfill
    \begin{subfigure}[b]{0.48\textwidth}
        \centering
        \includegraphics[width=\textwidth]{25_itu_attenuation_breakdown_bar.png}
        \caption{Attenuation breakdown at $99.9\%$ availability.}
        \label{fig:itu_breakdown_bar}
    \end{subfigure}
    \caption{ITU-R P.618 slant-path propagation impairments and specific contribution analysis for the Rome-Ankara link.}
    \label{fig:itu_propagation_metrics}
\end{figure}

\subsection[Monte-Carlo Rain Outage Simulation]{Monte-Carlo Rain Outage Simulation}

\subsubsection{Weather-State Model}
The ITU-R design point gives attenuation exceeded for a specified percentage of an average year; it does not provide a chronological weather sequence. A separate Monte Carlo model is therefore used to generate an illustrative distribution of hourly margins.

\subsubsection{Availability Estimation Rationale}
This stochastic model is a sensitivity exercise based on assumed state probabilities and fade ranges. It is not fitted to measured precipitation records.

\subsubsection{Mathematical Formulation}
The simulation executes $N=8760$ independent hourly trials. At each sample $t_i$, a weather state $S_i$ is drawn from a probability mass function:
\begin{itemize}
    \item \textbf{$S_{\mathrm{clear}}$:} A small residual excess attenuation uniformly sampled between 0 and 0.15~dB.
    \item \textbf{$S_{\mathrm{light}}$ and $S_{\mathrm{moderate}}$:} Small to medium fade depths drawn from uniform attenuation ranges.
    \item \textbf{$S_{\mathrm{heavy}}$:} Rare high-fade states that represent severe Ku-band rain attenuation.
\end{itemize}
The discrete state variable is therefore:
\begin{equation}
    S_i \in \{S_{\mathrm{clear}}, S_{\mathrm{light}}, S_{\mathrm{moderate}}, S_{\mathrm{heavy}}\}.
\end{equation}
For each generated hourly weather state, the active link margin is recalculated as:
\begin{equation}
    M_i =
    \left(\frac{E_b}{N_0}\right)_i
    -
    \left(\frac{E_b}{N_0}\right)_{\mathrm{req}}.
\end{equation}
The outage indicator for each sample is:
\begin{equation}
    \mathcal{O}_i =
    \begin{cases}
    1, & M_i < 0, \\
    0, & M_i \ge 0.
    \end{cases}
\end{equation}
The corresponding link availability is then estimated as:
\begin{equation}
    A_{\mathrm{link}} =
    100
    \left(
    1 -
    \frac{1}{N}
    \sum_{i=1}^{N}
    \mathcal{O}_i
    \right)\%.
\end{equation}

\subsubsection{Operational Impact and Results}
The output is a distribution of link margin under the assumed states and random seed. Because successive samples are independent, it does not reproduce storm duration or temporal correlation.

\textbf{Note:} This Monte Carlo model is an illustrative stochastic availability simulation based on assumed weather-state probabilities (light rain: 2.0\%, moderate rain: 0.6\%, heavy rain: 0.1\%), not a site-calibrated meteorological time-series model derived from measured annual precipitation data. Results should therefore be interpreted as \textit{estimated availability under the assumed stochastic weather-state model}, not as a prediction of actual annual availability.

\begin{figure}[H]
    \centering
    \includegraphics[width=0.75\textwidth]{16_availability_margin_histogram.png}
    \caption{Probability density distribution of the dynamic link margin over $8{,}760$ simulated hours of the meteorological year (Monte Carlo stochastic analysis).}
    \label{fig:mc_availability_histogram}
\end{figure}

\subsection[DVB-S2 ACM (Adaptive Coding \& Modulation)]{DVB-S2 ACM (Adaptive Coding \& Modulation)}

\subsubsection{Adaptive-Throughput Context}
A fixed MODCOD must be selected to meet its target availability with adequate margin. When channel conditions improve, that fixed operating point may leave capacity unused.

\subsubsection{Fixed-Rate and ACM Metrics}
The fixed-rate link margin and the ACM throughput analysis are treated as separate performance views. The former evaluates closure against a prescribed $(E_b/N_0)$ threshold, whereas the latter maps the available $(E_s/N_0)$ to DVB-S2 MODCOD thresholds and net throughput for a specified carrier bandwidth and roll-off factor. The two metrics should not be interpreted as the same quantity.

ACM maps the measured $E_s/N_0$ to a modulation and code rate. Higher-efficiency modes are selected when margin is available; more robust modes are selected during fades.

\subsubsection{Mathematical Formulation}
The receiver continuously monitors $(E_s/N_0)_{\mathrm{available}}$ and dynamically selects the highest spectral efficiency $\eta$ (in bits/symbol) from the DVB-S2 \cite{dvbs2} MODCOD ladder that satisfies:
\begin{equation}
    (E_s/N_0)_{\mathrm{available}} \ge (E_s/N_0)_{\mathrm{MODCOD,req}} + M_{\mathrm{impl}}
\end{equation}

\begin{figure}[H]
    \centering
    \includegraphics[width=0.75\textwidth]{28_dvbs2_modcod_ladder.png}
    \caption{DVB-S2 MODCOD ladder: trade-off between spectral efficiency and required signal-to-noise ratio.}
    \label{fig:dvbs2_modcod_ladder}
\end{figure}

Using the usable symbol rate $R_{\mathrm{s}} = B / (1 + \alpha_{\mathrm{rolloff}})$, the active net throughput is:
\begin{equation}
    R_{\mathrm{net}}\,[\mathrm{Mbps}] = \eta \cdot R_{\mathrm{s}}
\end{equation}

\subsubsection{Operational Impact and Results}
For a 36~MHz carrier with $\alpha=0.20$, the symbol rate is 30~Msps. The implemented threshold table selects QPSK 3/4 for the clear-sky thermal and interference cases and QPSK 1/2 at the 99.9\% atmospheric design point.

Table~\ref{tab:acm_consistency} lists the available and required $E_s/N_0$ values used by the selector.

\begin{table}[H]
    \centering
    \renewcommand{\arraystretch}{1.2}
    \caption{DVB-S2 ACM Consistency and Operating Point Verification}
    \label{tab:acm_consistency}
    \vspace{0.1cm}
    \resizebox{\textwidth}{!}{
      \begin{tabular}{l c c c c c}
        \toprule
        \textbf{Case} & \textbf{Avail. $E_s/N_0$} & \textbf{Selected MODCOD} & \textbf{Req. $E_s/N_0$} & \textbf{Efficiency ($\eta$)} & \textbf{Net Throughput ($\eta R_{\mathrm{s}}$)} \\
        \midrule
        \textbf{Clear-Sky (Thermal)} & $\approx 6.38$ dB & QPSK 3/4 & 4.03 dB & 1.487 & $\approx 44.61$ Mbps \\
        \textbf{With ASI \& IMD} & $\approx 5.92$ dB & QPSK 3/4 & 4.03 dB & 1.487 & $\approx 44.61$ Mbps \\
        \textbf{99.9\% Rain Fade} & $\approx 2.20$ dB & QPSK 1/2 & 1.00 dB & 0.989 & $\approx 29.67$ Mbps \\
        \bottomrule
    \end{tabular}
    }
    \vspace{0.1cm}
    \small\textit{Note: All points assume 1.0 dB implementation margin. Throughput is computed using $R_{\mathrm{s}} = 30$ Msps.}\normalsize
\end{table}

\begin{figure}[H]
    \centering
    \includegraphics[width=0.75\textwidth]{29_dvbs2_acm_vs_availability.png}
    \caption{DVB-S2 ACM dynamic throughput and MODCOD selection over the link availability spectrum.}
    \label{fig:dvbs2_acm_availability}
\end{figure}

\subsection[Uplink Power Control (ULPC)]{Uplink Power Control (ULPC)}

Uplink power control is discussed as a mitigation option but is not implemented as a closed-loop simulation. An ideal first-order controller would increase earth-station transmit power by the estimated uplink fade:
\begin{equation}
    P_{\mathrm{tx,dynamic}}\,[\mathrm{dBW}] = P_{\mathrm{tx,nominal}}\,[\mathrm{dBW}] + A_{\mathrm{rain,up}}\,[\mathrm{dB}]
\end{equation}
subject to $P_{\mathrm{tx,dynamic}} \le P_{\mathrm{tx,max}}$. Available HPA headroom, EIRP limits, ASI constraints, transponder loading, and control-loop delay limit the recoverable fade. No ULPC gain is credited in the reported margins.
